\documentclass[main.tex]{subfiles}


\begin{document}

%from pagenumber to up
% \phantomsection 
% \hypertarget{toc}{}

 

 

 
   

\section[Закон Да́льтона]{Закон Д\'альтона}

\sbox{0}{%
    \begin{tikzpicture}[font=\small,            line cap=round,                             >={Stealth[inset=0pt,round,length=3mm, angle'=20]},vec/.style={>={Stealth[inset=0pt,round,length=3.5mm, angle'=20]}}]


%gas
\draw (0,0) rectangle++ (3,3);

% buff and radius
\def\bu{2.84}
\def\ra{.04}

\foreach \i in {1,...,50}
  \fill[Green,shift={(.08,0.08)}] ({rnd*\bu cm}, {rnd*\bu cm}) circle (\ra);

\foreach \i in {1,...,50}
  \fill[red,shift={(.08,0.08)}] ({rnd*\bu cm}, {rnd*\bu cm}) circle (\ra);

    \end{tikzpicture}%
}

{%
\setlength\intextsep{0pt}
\begin{wrapfigure}[10]{r}{\wd0}
    \centering
    \usebox{0}%
    \caption{Два газа в сосуде}
    \label{fig:p69}
\end{wrapfigure}

В природе и в технике очень часто имеют дело со \emph{смесью} нескольких газов\footnote{Везде далее газы считаются идеальными.}. Например, воздух~--- это смесь азота, кислорода, аргона, углекислого газа и других газов.


Пусть имеется смесь двух газов в сосуде (рис.\,\ref{fig:p69}).

%     \begin{figure}[H]
%     \centering

%     \begin{tikzpicture}[font=\small,            line cap=round,                             >={Stealth[inset=0pt,round,length=3mm, angle'=20]},vec/.style={>={Stealth[inset=0pt,round,length=3.5mm, angle'=20]}}]


% %gas
% \draw (0,0) rectangle++ (3,3);

% % buff and radius
% \def\bu{2.86}
% \def\ra{.04}

% \foreach \i in {1,...,50}
%   \fill[Green,shift={(.08,0.08)}] ({rnd*\bu cm}, {rnd*\bu cm}) circle (\ra);

% \foreach \i in {1,...,50}
%   \fill[red,shift={(.08,0.08)}] ({rnd*\bu cm}, {rnd*\bu cm}) circle (\ra);

%     \end{tikzpicture}
% \caption{Два газа в сосуде}
% \label{fig:p69}
% \end{figure}

Газы на рис.\,\ref{fig:p69} имеют одинаковую температуру. Молекулы газов для наглядности окрашены~--- можно говорить как бы о зеленом и красном газах. Их смесь оказывает на стенки давление, обозначаемое~$P_\text{зк}$. В любой смеси газов каждый из газов ведет себя независимо от других газов, то есть зеленый и красный газы в предложенном примере можно рассматривать по отдельности в данном сосуде (рис.\,\ref{fig:p70}).  

}

    \begin{figure}[H]
    \centering

    \begin{tikzpicture}[font=\small,            line cap=round,                             >={Stealth[inset=0pt,round,length=3mm, angle'=20]},vec/.style={>={Stealth[inset=0pt,round,length=3.5mm, angle'=20]}}]


%gas
\draw (0,0) rectangle++ (3,3);

% buff and radius
\def\bu{2.84}
\def\ra{.04}

\foreach \i in {1,...,50}
  \fill[Green,shift={(.08,0.08)}] ({rnd*\bu cm}, {rnd*\bu cm}) circle (\ra);

%%%%%%%%%%%%%%%%%%%%%%%%%%%%%

\begin{scope}[xshift=5cm]
    
%gas
\draw (0,0) rectangle++ (3,3);

% buff and radius
\def\bu{2.84}
\def\ra{.04}

\foreach \i in {1,...,50}
  \fill[red,shift={(.08,0.08)}] ({rnd*\bu cm}, {rnd*\bu cm}) circle (\ra);
  
\end{scope}


    \end{tikzpicture}
\caption{Газы по отдельности}
\label{fig:p70}
\end{figure}

На рис.\,\ref{fig:p70} (слева) из сосуда удалены все газы кроме зеленого; оставшийся газ производит давление, обозначаемое~$P_\text{з.\,см}$ (подпись~<<см>> напоминает, что у данного газа такое же состояние, какое он имеет в смеси). Аналогично, на рис.\,\ref{fig:p70} (справа) в сосуде оставили только красный газ, производящий давление~$P_\text{к.\,см}$.

\textbf{Парциальное давление}~--- это давление,  которое производил бы выбранный газ, входящий в состав смеси, если удалить остальные газы из сосуда. Таким образом, в рассматриваемой ситуации с зеленым и красным газами их давления~$P_\text{з.\,см}$ и~$P_\text{к.\,см}$~--- это парциальные давления.
\begin{law}
\textbf{Закон Д\'альтона.} Давление смеси газов есть
сумма их парциальных давлений:
\begin{equation}
    \label{5dal_e1}
    P_\text{см}=P_\text{1\,см}+P_\text{2\,см}+\ldots,
\end{equation}
где $P_\text{1\,см},P_\text{2\,см},\ldots$~--- парциальные давления первого, второго и так далее газов.  
\end{law}
Так, для ситуации, проиллюстрированной рисунками~\ref{fig:p69} и~\ref{fig:p70}, уравнение~\eqref{5dal_e1} дает: $P_\text{зк}=P_\text{з.\,см}+P_\text{к.\,см}$.

Теперь можно разобрать стандартную задачу на смеси газов.

\medskip

\noindent\textbf{Задача.} Соединенные краном сосуды с газами под давлением 100 и 600~кПа имеют объемы 2~л и 3~л соответственно. Какое установится давление, если кран открыть? Температура постоянна.

\smallskip

\noindent\emph{Решение.} Парциальные давления дает уравнение Менделеева"--~Клапейрона:\linebreak
$P_\text{1\,см}=\dfrac{\nu_1 RT}{V_\text{см}}=\dfrac{P_1V_1}{V_\text{см}}=40$~кПа
% \dfrac{100\cdot10^{3}\cdot2\cdot10^{-3}}{5\cdot10^{-3}}=40$~кПа
и 
$P_\text{2\,см}=\dfrac{\nu_2 RT}{V_\text{см}}=\dfrac{P_2V_2\strut}{V_\text{см}\strut}=360$~кПа. Давление смеси по закону Д\'альтона равно: $P_\text{см}=P_\text{1\,см}+P_\text{2\,см}=400\strut$~кПа.







 
\end{document}